% Part: first-order-logic
% Chapter: tableaux
% Section: provability-consistency

\documentclass[../../../include/open-logic-section]{subfiles}

\begin{document}

\iftag{FOL}
      {\olfileid[zh]{fol}{tab}{prv}}
      {\olfileid[zh]{pl}{tab}{prv}}

\olsection{\usetoken{S}{derivability}与一致性}

我们现在要确立!!{derivability}关系的若干性质。它们本身彼此独立地具有意义，但它们每一个都将在完全性定理的证明中发挥作用。

\begin{prop}\ollabel{prop:provability-contr}
  如果 $\Gamma \Proves !A$ 且 $\Gamma \cup \{!A\}$ 是不一致的，则 $\Gamma$ 是不一致的。
\end{prop}

\begin{proof}
存在有穷的 $\Gamma_0 = \{!B_1, \dots, !B_n\}$ 和 $\Gamma_1
  =\{!C_1, \dots, !C_n\} \subseteq \Gamma$，使得
  \begin{align*}
    \{\sFmla{\False}{!A}, &
    \sFmla{\True}{!B_1}, \dots, \sFmla{\True}{!B_n}\} \\
    \{\sFmla{\True}{!A}, &
    \sFmla{\True}{!C_1}, \dots, \sFmla{\True}{!C_m}\}
  \end{align*}
  具有闭的!!{tableau}。利用作用在 $!A$ 上的\Cut{}规则，我们可以把这些组合成单一的闭!!{tableau}，它表明 $\Gamma_0
  \cup \Gamma_1$ 是不一致的。由于 $\Gamma_0
  \subseteq \Gamma$ 且 $\Gamma_1 \subseteq \Gamma$，故 $\Gamma_0 \cup
  \Gamma_1 \subseteq \Gamma$，因此 $\Gamma$~是不一致的。
\end{proof}

\begin{prop}
\ollabel{prop:prov-incons}
$\Gamma \Proves !A$ 当且仅当 $\Gamma \cup \{\lnot !A\}$ 是不一致的。
\end{prop}

\begin{proof}
首先设 $\Gamma \Proves !A$，即存在
\[
\{\sFmla{\False}{!A},
\sFmla{\True}{!B_1}, \dots, \sFmla{\True}{!B_n}\}
\]
的一个闭!!{tableau}。利用$\TRule{\True}{\lnot}$规则，这可以转变为
\[
\{\sFmla{\True}{\lnot !A},
\sFmla{\True}{!B_1}, \dots, \sFmla{\True}{!B_n}\} \text{。}
\]
的一个闭!!{tableau}。

另一方面，如果后者的闭!!{tableau}存在，我们可以去掉每一个由\TRule{\True}{\lnot}作用于第一个假设~$\sFmla{\True}{\lnot !A}$而得到的公式以及该假设本身，并添加假设 $\sFmla{\False}{!A}$，从而把它转变为前者的闭!!{tableau}。因为如果某条支之前由于包含 \TRule{\True}{\lnot} 作用于 $\sFmla{\True}{\lnot !A}$ 的结论（即 $\sFmla{\False}{!A}$）而是闭的，那么新!!{tableau}中对应的支也是闭的。如果旧表列中的某条支由于同时包含假设 $\sFmla{\True}{\lnot !A}$ 与 $\sFmla{\False}{\lnot
  !A}$ 而是闭的，我们可以通过把\TRule{\False}{\lnot} 应用于 $\sFmla{\False}{\lnot !A}$ 以得到
$\sFmla{\True}{!A}$，从而把它转变为闭支。由于我们添加了
$\sFmla{\False}{!A}$ 作为假设，这使该支闭合。
\end{proof}

\begin{prob}
证明：$\Gamma \Proves \lnot !A$ 当且仅当 $\Gamma \cup \{!A\}$ 是不一致的。
\end{prob}

\begin{prop}\ollabel{prop:explicit-inc}
  如果 $\Gamma \Proves !A$ 且 $\lnot !A \in \Gamma$，则 $\Gamma$ 是
  不一致的。
\end{prop}

\begin{proof}
  设 $\Gamma \Proves !A$ 且 $\lnot !A \in \Gamma$。则存在
  $!B_1$, \dots, $!B_n \in \Gamma$，使得 \
  \[
  \{\sFmla{\False}{!A}, \sFmla{\True}{!B_1}, \dots, \sFmla{\True}{!B_n}\}
  \]
  具有闭的表列。把假设 \sFmla{\False}{!A} 替换为
  \sFmla{\True}{\lnot !A}，并把这些假设之后插入
  \TRule{\True}{\lnot} 作用于 \sFmla{\False}{!A} 的结论。在!!{tableau}中通过诉诸旧!!{tableau}中的第~$1$ 行而得到辩护的任何!!{sentence}，现在都通过诉诸第~$n+1$ 行而得到辩护。因此，如果旧!!{tableau}是闭的，新表列也是闭的。它表明 $\Gamma$ 是不一致的，因为所有假设都在~$\Gamma$ 中。
\end{proof}

\begin{prop}\ollabel{prop:provability-exhaustive}
  如果 $\Gamma \cup \{!A\}$ 与 $\Gamma \cup \{\lnot !A\}$ 两者都是
  不一致的，则 $\Gamma$ 是不一致的。
\end{prop}

\begin{proof}
  如果存在 $!B_1$, \dots, $!B_n \in \Gamma$ 以及 $!C_1$, \dots,
  $!C_m \in \Gamma$，使得
  \begin{align*}
    \{\sFmla{\True}{!A}, &
    \sFmla{\True}{!B_1}, \dots, \sFmla{\True}{!B_n}\} \text{ 以及}\\
    \{\sFmla{\True}{\lnot !A}, &
    \sFmla{\True}{!C_1}, \dots, \sFmla{\True}{!C_m}\}
  \end{align*}
  两者都具有闭!!{tableau}，那么我们可以通过以下方式构造单一的、组合的!!{tableau}，它表明 $\Gamma$ 是不一致的：以
  $\sFmla{\True}{!B_1}$, \dots, $\sFmla{\True}{!B_n}$
  连同 $\sFmla{\True}{!C_1}$, \dots, $\sFmla{\True}{!C_m}$ 作为
  假设，然后应用一次\Cut{}规则。这产生两条
  支，一条以 $\sFmla{\True}{!A}$ 开头，另一条以
  $\sFmla{\False}{!A}$ 开头。

  在左侧，添加第一个!!{tableau}的假设以下的部分。在这里，每一种规则应用仍然
  是正确的，因为第一个!!{tableau}的每一个假设（包括
  $\sFmla{\True}{!A}$）都是可用的。因此，$\sFmla{\True}{!A}$ 之下的每一条支
  都闭合。

  在右侧，添加第二个!!{tableau}的假设以下的部分，去掉
  把~$\TRule{\True}{\lnot}$ 应用于 $\sFmla{\True}{\lnot !A}$ 的任何应用结果。把 $\TRule{\True}{\lnot}$ 应用于 $\sFmla{\True}{\lnot !A}$ 的结论是
  $\sFmla{\False}{!A}$，它毕竟仍是可用的，因为它是组合!!{tableau}右侧\Cut{}规则的结论。

  如果第二个表列中的某条支由于包含
  假设 $\sFmla{\True}{\lnot !A}$（它在组合!!{tableau}中不再作为
  假设出现）以及
  $\sFmla{\False}{\lnot !A}$ 而是闭的，我们可以把\TRule{\False}{\lnot}
  应用于 $\sFmla{\False}{\lnot !A}$ 以得到 $\sFmla{\True}{!A}$。现在组合!!{tableau}中对应的支也闭合，
  因为它包含\Cut{}规则的右侧结论
  $\sFmla{\False}{!A}$。如果第二个!!{tableau}中的某条支由于任何其他原因闭合，那么组合!!{tableau}中对应的支也闭合，因为除 $\sFmla{\True}{\lnot !A}$ 之外在旧第二个!!{tableau}的该支上出现的任何!!{signed formula}也出现在组合!!{tableau}的对应支上。
  \end{proof}

\end{document}
